\documentclass[a4paper,answers]{exam} % Pour la version avec les réponses, utiliser: \documentclass[answers]{exam}
\usepackage[utf8]{inputenc}
\usepackage[T1]{fontenc}
\usepackage{amsmath, amssymb}
\usepackage{siunitx}

\newcommand{\pts}[1]{\textbf{#1 pts}}

\extrawidth{1in}

%\firstpageheader{
%\large{Évaluation - Projet 2 : Modélisation moléculaire}

\pagestyle{headandfoot}
\runningheadrule
\firstpageheader{PHY 4268}{Évaluation Chapitre 2 - Méthodes d’optimisation géométrique}{28 Mars 2025}
\runningheader{PHY 4268}
{Évaluation Chapitre 2 - Méthodes d’optimisation géométrique, Page \thepage\ of \numpages}
{28 Mars 2025}
\firstpagefooter{}{}{}
\runningfooter{}{}{}

\begin{document}

\begin{itemize}
    \item Durée : 20 minutes
    \item Type d'épreuve : Questions à réponses courtes et questions de compréhension
\end{itemize}

\vspace{0.2cm}

\begin{questions}

\question[5] \textbf{Objectif de l'optimisation géométrique.} En termes simples, quel est le \textbf{but principal} de l'optimisation géométrique en chimie quantique computationnelle ? (Maximum 2 phrases)

\begin{solution}
\begin{itemize}
    \item[] But principal : Trouver la géométrie moléculaire correspondant à l'énergie la plus basse / structure moléculaire stable \hfill \marks[3]
    \item[] Précision/Complément :  Ajuster les positions des atomes pour minimiser l'énergie totale / trouver la configuration électronique la plus stable \hfill \marks[2]
\end{itemize}
\end{solution}


\question[5] \textbf{Surface d'Énergie Potentielle (PES) et Minima.}
\begin{parts}
    \part[2] Définissez brièvement ce qu'est la \textbf{Surface d'Énergie Potentielle (PES)} dans le contexte de la modélisation moléculaire. (1 phrase)
    \begin{solution}
    \begin{itemize}
        \item[] Définition PES : Surface représentant l'énergie moléculaire en fonction des positions des noyaux \hfill \marks[2]
    \end{itemize}
    \end{solution}

    \part[3] Quel élément de la PES recherche-t-on lors d'une optimisation géométrique, et à quoi correspond cet élément en termes de structure moléculaire ? (2 phrases maximum)
    \begin{solution}
    \begin{itemize}
        \item[] Élément recherché : Minimum de la PES \hfill \marks[1]
        \item[] Correspondance : Correspond à la géométrie d'équilibre / structure stable de la molécule \hfill \marks[2]
    \end{itemize}
    \end{solution}
\end{parts}


\question[5] \textbf{Rôle et Avantage de l'Hessien.}
\begin{parts}
    \part[3] Quel est l'\textbf{avantage principal} d'utiliser le \textbf{Hessien} (matrice des dérivées secondes de l'énergie) dans un algorithme d'optimisation géométrique, par rapport à un algorithme utilisant uniquement le gradient ? (1-2 phrases maximum)
    \begin{solution}
    \begin{itemize}
        \item[] Avantage principal de l'Hessien : Convergence plus rapide vers le minimum \hfill \marks[2]
        \begin{itemize}
            \item[] Précision/Complément : L'Hessien donne des informations sur la courbure de la PES, permettant des pas d'optimisation plus efficaces / prédit mieux la direction du minimum \hfill \marks[1]
        \end{itemize}
    \end{itemize}
    \end{solution}

    \part[2] Le calcul du Hessien est-il toujours effectué lors d'une optimisation géométrique ? \textbf{Justifiez brièvement} votre réponse. (1 phrase)
    \begin{solution}
    \begin{itemize}
        \item[] Calcul du Hessien toujours effectué : Non \hfill \marks[1]
        \begin{itemize}
            \item[] Justification : Calcul du Hessien coûteux en temps de calcul / algorithmes comme descente de gradient n'utilisent pas le Hessien / compromis coût-efficacité \hfill \marks[1]
        \end{itemize}
    \end{itemize}
    \end{solution}
\end{parts}


\question[5] \textbf{Critères de Convergence.}
\begin{parts}
    \part[2] Citez \textbf{deux critères de convergence} typiquement utilisés pour arrêter une optimisation géométrique. (2 critères)
    \begin{solution}
    \textbf{Corrigé Indicatif et Barème Détaillé (Question 4a - 2 points)}
    \begin{itemize}
        \item[] Critère 1 : Gradient de l'énergie (force) \hfill \marks[1]
        \item[] Critère 2 : Déplacement atomique OU Convergence de l'énergie \hfill \marks[1]
        \item[] (Accepter d'autres critères pertinents mentionnés dans le cours)
    \end{itemize}
    \end{solution}
    \fbox{\parbox{0.9\textwidth}{\vspace{0.5cm}}} % Boîte pour la réponse étudiant

    \part[3] L'optimisation géométrique garantit-elle toujours de trouver le \textbf{minimum global} de la PES ?  \textbf{Justifiez brièvement} votre réponse. (2 phrases maximum)
    \begin{solution}
    \textbf{Corrigé Indicatif et Barème Détaillé (Question 4b - 3 points)}
    \begin{itemize}
        \item[] Garantie minimum global : Non \hfill \marks[1]
        \begin{itemize}
            \item[] Justification : Optimisation converge souvent vers un minimum local proche de la géométrie initiale, et non nécessairement vers le minimum global \hfill \marks[2]
        \end{itemize}
    \end{itemize}
    \end{solution}
\end{parts}


\end{questions}

\end{document}
